
\Problem{Permutacions quasi ordenades}

\Statement

Donat un enter $k\geq 0$, diem que una permutació d'$n$ elements 
$\{1,\dots,n\}$ està quasi ordenada si com a màxim hi ha $k$ elements
que estan en una posició diferent de la seva posició natural.
\medskip

Per exemple, aquestes són totes les permutacions quasi ordenades per a $n=6$
i $k=0$ (i també per a $k=1$):

\begin{verbatim}
1 2 3 4 5 6
\end{verbatim}

i aquestes són totes les permutacions quasi ordenades per a $n=5$ i $k=2$:

\begin{verbatim}
1 2 3 4 5
1 2 3 5 4
1 2 4 3 5
1 2 5 4 3
1 3 2 4 5
1 4 3 2 5
1 5 3 4 2
2 1 3 4 5
3 2 1 4 5
4 2 3 1 5
5 2 3 4 1
\end{verbatim}

\medskip

Feu un programa de generació exhaustiva que llegeixi dos nombres naturals
$n > 0$ i $k\geq 0$ i escrigui totes les permutacions quasi ordenades 
de $\{1,\dots,n\}$ en ordre lexicogràfic de tal manera que cada permutació
tingui com a màxim $k$ elements en una posició diferent de la seva posició 
natural.

\Sample
